\subsection{Applications : I. Canonical decompositions of rational numbers and of rational functions in one indeterminate}

THEOREM 2. --- Let A be a principal ideal domain, let K be its field of fractions and let P be a system of representatives of irreducible elements of A. Given an element \(x \in K\) , there exist a finite subset H of P, elements \(a_{0} \in A\) and \(a_{p} \in A\) not divisible by p in A ( \(p \in H\) ), and integers \(s(p) > 0\) ( \(p \in H\) ) such that

\[
x = a _ {0} + \sum_ {p \in \mathrm{H}} a _ {p} p ^ {- s (p)},\tag{1}
\]

where H and the \(s(p)\) are uniquely determined by these conditions.

Moreover, if \(R_{p}\) denotes a subset of A containing precisely one element of each residue class mod p in A ( \(p \in P\) ), then each \(x \in K\) can be uniquely expressed in the form

\[
x = a + \sum_ {p \in P} \left(\sum_ {h = 1} ^ {\infty} r _ {p h} p ^ {- h}\right)\tag{2}
\]

where \(a \in \mathbf{A}\) and \(r_{ph} \in \mathbf{R}_p\) for all \(h\) and \(p\) , and all but finitely many of the \(r_{ph}\) are zero.

Consider K as an A-module; then A is the submodule generated by 1. The quotient module K/A is the quotient of K by the equivalence relation \(x' - x \in \mathbf{A}\) , which is also written, in the notation of VI, p.~6, as \(x \equiv x' \pmod{1}\) ; let \(f\) denote the natural homomorphism from K onto M = K/A.

The quotient module M is a torsion module, for every element of M has the form \(f(a/b)\) ( \(a \in A\) , \(b \in A\) , \(b \neq 0\) ), whence \(bf(a/b) = f(a) = 0\) . Hence Th. 1 of VII, p.~8 applies. Let \(M_{p}\) ( \(p \in P\) ) be the submodule of elements of M annihilated by powers of p; then \(f^{-1}(M_{p})\) is the subring \(A_{p}\) of elements of K of the form \(ap^{-n}\) where \(a \in A\) and \(n \geqslant 0\) is an integer. The module M is the direct sum of the \(M_{p}\) , so every \(x \in K\) is congruent mod 1 to an element in the sum of the \(A_{p}\) ; in other words, formula (1) holds, with the \(s(p)\) integers \textgreater0, and the \(a_{p}\) finitely many elements of A such that \(a_{p}\) is not a multiple of p.

We now show that these conditions on the \(s(p)\) and the \(a_{p}\) completely determine H and the \(s(p)\) . Indeed H is then the set of \(p \in \mathbf{P}\) such that the component of \(f(x)\) in \(\mathbf{M}_{p}\) is \(\neq 0\) . On the other hand if \(s\) and \(s'\) are two integers \(>0\) such that \(s \geqslant s'\) and if \(a\) and \(a'\) are elements of \(A\) not divisible by \(p\) such that \(ap^{-s} \equiv a'p^{-s'} \pmod{1}\) , then we deduce that \(a \equiv a'p^{s - s'} \pmod{p^s}\) ; if \(s > s'\) then \(a \equiv 0 \pmod{p}\) , contradicting the hypotheses. This argument also shows that each \(a_p\) is well defined mod \(p^{s(p)}\) .

To complete the proof, we note first of all that in each residue class mod \(p^s\) in \(\mathbf{A}\) there exists a unique element of the form \(\sum_{h=0}^{s-1} r_h p^h\) with \(r_h \in \mathbb{R}_p\) for \(0 \leqslant h \leqslant s-1\) . In fact, we proceed by induction on \(s\) (the property follows from the definition of \(\mathbb{R}_p\) for \(s=1\) ): let \(x \in \mathbf{A}\) ; by hypothesis there is a unique element of the form \(\sum_{h=0}^{s-2} r_h p^h\) ( \(r_h \in \mathbb{R}_p\) ) in the residue class of \(x \mod p^{s-1}\) ; then \(x - \sum_{h=0}^{s-2} r_h p^h\) is a multiple \(ap^{s-1}\) of \(p^{s-1}\) ; now there exists a unique element \(r_{s-1}\) of \(\mathbb{R}_p\) such that \(a \equiv r_{s-1} (\bmod p)\) ; whence \(x \equiv \sum_{h=0}^{s-1} r_h p^h (\bmod p^s)\) . To obtain formula (2) it is now enough to apply this fact to each \(a_p\) in formula (1). The uniqueness is clear in view of the above.

The following are the most important applications of Th. 2:

I. The ring A is the ring Z of rational integers, and K = Q. Let P be the set of prime numbers \textgreater{} 0, and for each \(p \in P\) let \(R_{p}\) be the interval \((0, p - 1)\) in Z. Thus we have the canonical decomposition

\[
x = a + \sum_ {p \in P} \left(\sum_ {h = 1} ^ {\infty} e _ {p h} p ^ {- h}\right)
\]

with \(a \in \mathbf{Z}\) , \(e_{ph} \in \mathbf{Z}\) and \(0 \leqslant e_{ph} \leqslant p - 1\) .

\begin{enumerate}
\def\labelenumi{\Roman{enumi}.}
\setcounter{enumi}{1}
\tightlist
\item
  The ring A is the ring E{[}X{]} of polynomials in one indeterminate over a commutative field E, and \(\mathbf{K} = \mathbf{E}(\mathbf{X})\) . Let P be the set of monic irreducible polynomials in E{[}X{]} (VII, p.~5). For \(p \in P\) we can, by virtue of euclidean division of polynomials (IV, p.~10), take \(R_{p}\) to be the set of polynomials of degree strictly less than that of p.~Thus we have the decomposition (called canonical) of a rational function \(r(\mathbf{X}) \in \mathbf{E}(\mathbf{X})\) :
\end{enumerate}

\[
r (\mathrm{X}) = a (\mathrm{X}) + \sum_ {p \in \mathrm{P}} \left(\sum_ {h = 1} ^ {\infty} v _ {p h} (\mathrm{X}) \cdot p (\mathrm{X}) ^ {- h}\right)
\]

where \(a(\mathbf{X})\) is a polynomial and \(v_{ph}(\mathbf{X})\) is a polynomial of degree strictly less than that of \(p(\mathbf{X})\) , for all \(p\) and \(h\) . In particular, if the field \(\mathbf{E}\) is algebraically closed, then the \(p(\mathbf{X})\) have the form \(\mathbf{X} - \alpha\) with \(\alpha \in \mathbf{E}\) , and the \(v_{ph}(\mathbf{X})\) are thus constants.

Hence we can say that the vector space \(E(X)\) over the field E has a basis consisting of the monomials \(X^{n}\) ( \(n \geqslant 0\) an integer) and the rational functions of the form \(\mathbf{X}^{m}/(p(\mathbf{X}))^{h}\) , where \(p \in P\) and h and m are integers with \(h \geqslant 1\) and \(0 \leqslant m < \deg(p)\) .

